\input{sol_pre.tex}
\input{q_2026-03-02_2.tex}
\begin{document}
\input{sol_head.tex}

%================= 1
\Savol{1}
$2026=7\cdot289+3\Rightarrow 2026^{2025}\equiv3^{2025}\pmod7$. \ $3^6\equiv1\pmod7$ (Ferma), \ $2025=6\cdot337+3$:
\[3^{2025}\equiv3^3=27\equiv6\pmod7.\]
\Javob{C) 6}

%================= 2
\Savol{2}
1) $\frac7{17}$, 2) $\frac{22}{17}$ — ratsional. \ 3) $\lg3$ irratsional $\Rightarrow$ $2+\lg3$ irratsional.\\
4) $\sqrt{\ln e^4}+\sqrt{\sqrt[3]8-1}=\sqrt4+\sqrt1=3$ — ratsional. \ 5) $\sqrt[3]{1+\sqrt2}$ — irratsional.
\Javob{B) 3 ta}

%================= 3
\Savol{3}
Birlar raqami $b$, önlar $b+3$: son $11b+30$, raqamlar yiğindisi $2b+3$.
\[(2b+3)(11b+30)=364:\quad b=2\ \Rightarrow\ 7\cdot52=364.\]
Son $52$, raqamlar köpaytmasi $5\cdot2=10$.
\Javob{C) 10}

%================= 4
\Savol{4}
$x+y=205\,000$, \ $0{,}85x=1{,}2y\Rightarrow x=\dfrac{24}{17}y$ $\Rightarrow$ $\dfrac{41}{17}y=205\,000\Rightarrow y=85\,000$, \ $x=120\,000$.\\
Farq: $120\,000-85\,000=35\,000$.
\Javob{A) 35\,000}

%================= 5
\Savol{5}
\[\left(5^{\frac72\cdot\frac47}+16^{\frac3{12}}\right)\cdot3^{-3}=(25+2)\cdot\frac1{27}=1.\]
\Javob{C) 1}

%================= 6
\Savol{6}
$2\sqrt{11}=\sqrt{44}$, \ $6<\sqrt{44}<7$; \quad $7\sqrt5=\sqrt{245}$, \ $15<\sqrt{245}<16$.\\
Butun sonlar: $7,8,\dots,15$ — 9 ta.
\Javob{A) 9 ta}

%================= 7
\Savol{7}
\[\frac1{2+\sqrt2}-\frac1{2-\sqrt2}=\frac{(2-\sqrt2)-(2+\sqrt2)}{4-2}=-\sqrt2,\qquad -\sqrt2-\sqrt2=-2\sqrt2.\]
\Javob{B) $-2\sqrt2$}

%================= 8
\Savol{8}
Har qanday $n$ uçun $\dfrac{S_n}n=2n\Rightarrow S_n=2n^2$.
\[a_{10}=S_{10}-S_9=200-162=38.\]
\Javob{A) 38}

%================= 9
\Savol{9}
$b_1b_3=b_2^2$:
\[\frac1{b_1}+\frac1{b_2}+\frac1{b_3}=\frac{b_2b_3+b_1b_3+b_1b_2}{b_1b_2b_3}=\frac{b_2(b_1+b_2+b_3)}{b_2^3}=\frac{21}{b_2^2}=\frac7{12}\ \Rightarrow\ b_2^2=36,\ b_2=6.\]
\Javob{C) 6}

%================= 10
\Savol{10}
$u=a+2023$, \ $v=b+1$: \ $u+v=a+b+2024$.
\[(u+v)^2=4uv\ \Rightarrow\ (u-v)^2=0\ \Rightarrow\ a+2023=b+1\ \Rightarrow\ a-b=-2022.\]
\Javob{B) $-2022$}

%================= 11
\Savol{11}
$x\pm4\sqrt{x-4}=(x-4)\pm4\sqrt{x-4}+4=\left(\sqrt{x-4}\pm2\right)^2$; \ $x>8$ da $\sqrt{x-4}>2$:
\[\left(\sqrt{x-4}+2\right)+\left(\sqrt{x-4}-2\right)=2\sqrt{x-4}.\]
\Javob{A) $2\sqrt{x-4}$}

%================= 12
\Savol{12}
$1-\sin^4\alpha-\cos^4\alpha=1-\left(1-2\sin^2\alpha\cos^2\alpha\right)=2\sin^2\alpha\cos^2\alpha$:
\[\frac{2\sin^2\alpha\cos^2\alpha\cdot\cos^2\alpha}{\sin^2\alpha}=2\cos^4\alpha.\]
\Javob{A) $2\cos^4\alpha$}

%================= 13
\Savol{13}
Aniqlaniş sohasi $|x|\le2$. \ $\sqrt{4-x^2}=0$: $x=\pm2$; \ $\cos x=0$: $x=\pm\frac\pi2$ ($\frac\pi2<2$).\\
Jami 4 ta.
\Javob{B) 4 ta}

%================= 14
\Savol{14}
$\sqrt{7\pm4\sqrt3}=2\pm\sqrt3$, \ $(2+\sqrt3)(2-\sqrt3)=1$. \ $t=(2+\sqrt3)^x$:
\[t+\frac1t=14\ \Rightarrow\ t=7\pm4\sqrt3=(2\pm\sqrt3)^2\ \Rightarrow\ x=\pm2.\]
Yiğindi: $0$.
\Javob{D) 0}

%================= 15
\Savol{15}
$9-12x+4x^2=(2x-3)^2$, \ $6x^2-23x+21=(2x-3)(3x-7)$. Şart: $x>\frac73$, $x\ne\frac83$.\\
$A=\log_{3x-7}(2x-3)$: \ birinçi qöşiluvçi $2A$, ikkinçisi $1+\log_{2x-3}(3x-7)=1+\dfrac1A$.
\[2A+\frac1A+1=4\ \Rightarrow\ 2A^2-3A+1=0\ \Rightarrow\ A=1\ \text{yoki}\ A=\frac12.\]
$A=1$: $2x-3=3x-7\Rightarrow x=4$ — mos. \ $A=\frac12$: $(2x-3)^2=3x-7\Rightarrow4x^2-15x+16=0$, $D<0$.
\Javob{A) 4}

%================= 16
\Savol{16}
$t=|x|\ge0$: \ $t^2-(\pi-\sqrt{10})t-\sqrt{10}\pi=(t-\pi)(t+\sqrt{10})=0\Rightarrow t=\pi$.\\
$|x|=\pi\Rightarrow x=\pm\pi$ — 2 ta.
\Javob{A) 2}

%================= 17
\Savol{17}
Qaytma tenglama, $x^2$ ga bölamiz, $s=x+\frac1x$:
\[\left(x^2+\frac1{x^2}\right)-10\left(x+\frac1x\right)+26=0\ \Rightarrow\ s^2-10s+24=0\ \Rightarrow\ s=4\ \text{yoki}\ s=6.\]
$s=4$: $x=2\pm\sqrt3$; \ $s=6$: $x=3\pm2\sqrt2$. Eng kattasi $3+2\sqrt2$.
\Javob{D) $3+2\sqrt2$}

%================= 18
\Savol{18}
$x\ge\frac12$: $\bigl|x-1\bigr|\le7\Rightarrow\frac12\le x\le8$ — butunlar $1,\dots,8$ (8 ta).\\
$x<\frac12$: $|1-3x|\le7\Rightarrow-2\le x<\frac12$ — butunlar $-2,-1,0$ (3 ta). \ Jami 11 ta.
\Javob{A) 11}

%================= 19
\Savol{19}
$-1<1-\sqrt x<1\Rightarrow 0<\sqrt x<2\Rightarrow 0<x<4$. \ Butunlar: $1,2,3$.
\Javob{A) 3 ta}

%================= 20
\Savol{20}
$f(-5)=8-5=3$, \ $f(3)=2-3=-1$, \ $f(-1)=8-1=7$.
\Javob{B) 7}

%================= 21
\Savol{21}
$u=x^3-x+2$, \ $du=(3x^2-1)dx$; \ $u(0)=2$, \ $u(1)=2$:
\[\int_2^2u^{2025}\,du=0.\]
\Javob{A) 0}

%================= 22
\Savol{22}
$x_0=\frac\pi4$ da $\log_{\cos x}\sin x=1$, $\ln\frac\pi4<0$ $\Rightarrow$ modul içi manfiy, atrofda $f(x)=g(x)-\ln x$, $g(x)=\dfrac{\ln\sin x}{\ln\cos x}$.
\[g'(x)=\frac{\ctg x\cdot\ln\cos x+\tg x\cdot\ln\sin x}{\ln^2\cos x};\qquad x_0:\ \ln\sin x_0=\ln\cos x_0=-\tfrac12\ln2,\ \ g'(x_0)=\frac{2}{-\frac12\ln2}=-\frac4{\ln2}.\]
\[f'(x_0)=-\frac4{\ln2}-\frac1{x_0}=-\frac4\pi-\frac4{\ln2}.\]
\Javob{D) $-\frac4\pi-\frac4{\ln2}$}

%================= 23
\Savol{23}
$OB=OC$ $\Rightarrow$ $\angle OCB=40^\circ$, \ $\angle BOC=180^\circ-80^\circ=100^\circ$.\\
$\alpha=\angle BAC$ — içki çizilgan, $BC$ yoyga tiraladi: $\alpha=\frac12\angle BOC=50^\circ$.
\Javob{C) $50^\circ$}

%================= 24
\Savol{24}
$x^2=\dfrac{a^2-1}{(a-1)^2}=\dfrac{a+1}{a-1}$ $\Rightarrow$ $\dfrac{x^2+1}{x^2-1}=\dfrac{2a}{2}=a$.\\
$2y^2=\dfrac{1-b}{1+b}$ $\Rightarrow$ $\dfrac{2y^2-1}{2y^2+1}=\dfrac{-2b}{2}=-b$. \quad $4z^2=\dfrac{1-c^2}{(c-1)^2}=\dfrac{1+c}{1-c}$ $\Rightarrow$ $\dfrac{4z^2-1}{4z^2+1}=c$.
\[f=a-(-b)+c=a+b+c.\]
\Javob{A) $a+b+c$}

%================= 25
\Savol{25}
Bissektrisa xossasi: $\dfrac{AB}{BC}=\dfrac{AE}{EC}=\dfrac48\Rightarrow BC=12$.\\
Bissektrisa formulasi: $BE^2=AB\cdot BC-AE\cdot EC=72-32=40\Rightarrow BE=2\sqrt{10}$.
\Javob{A) $2\sqrt{10}$}

%================= 26
\Savol{26}
$AD:DB=2:1$, \ $BE:EC=1:1$. Menelay teoremasi ($\triangle ABE$, töğri çiziq $D$–$O$–$C$):
\[\frac{AD}{DB}\cdot\frac{BC}{CE}\cdot\frac{EO}{OA}=1\ \Rightarrow\ 2\cdot2\cdot\frac{EO}{OA}=1\ \Rightarrow\ AO:OE=4:1.\]
$S_{AEC}=\frac12S_{ABC}=30$, \ $S_{AOC}=\frac45\cdot30=24$.
\Javob{C) 24}

%================= 27
\Savol{27}
$CO$, $DO$ — $C$ va $D$ burçaklar bissektrisalari, $\angle C+\angle D=180^\circ\Rightarrow\angle COD=90^\circ$, \ $CD=\sqrt{4+16}=2\sqrt5$.\\
$r$ — $\triangle COD$ ning balandligi: $r=\dfrac{2\cdot4}{2\sqrt5}=\dfrac4{\sqrt5}$. Urinmalar: $C$ dan $\sqrt{4-\frac{16}5}=\frac2{\sqrt5}$, $D$ dan $\sqrt{16-\frac{16}5}=\frac8{\sqrt5}$.\\
$BC=r+\frac2{\sqrt5}=\frac6{\sqrt5}$, \ $AD=r+\frac8{\sqrt5}=\frac{12}{\sqrt5}$, \ balandlik $2r=\frac8{\sqrt5}$:
\[S=\frac{BC+AD}2\cdot2r=\frac9{\sqrt5}\cdot\frac8{\sqrt5}=\frac{72}5=14\frac25.\]
\Javob{D) $14\frac25$}

%================= 28
\Savol{28}
Qöşni tomonlar örtalari orasidagi masofa $a\cdot\frac{\sqrt3}2$, yuzlar har safar $\frac34$ marta kamayadi.
\[S=\frac{S_0}{1-\frac34}=4\cdot\frac{3\sqrt3}2=6\sqrt3.\]
\Javob{C) $6\sqrt3$}

%================= 29
\Savol{29}
$\overrightarrow{AB}\cdot\overrightarrow{AC}=24-24=0$ $\Rightarrow$ $\angle A=90^\circ$. \ $\overrightarrow{BC}=(-2;-14)$, \ $BC=\sqrt{200}=10\sqrt2$.\\
Töğri burçakli uçburçakda $R=\dfrac{BC}2=5\sqrt2$.
\Javob{C) $5\sqrt2$}

%================= 30
\Savol{30}
$HB=AH\ctg45^\circ=4$, \ $HC=AH\ctg60^\circ=\dfrac4{\sqrt3}$.
\[S_{HBC}=\frac12\cdot4\cdot\frac4{\sqrt3}\cdot\sin60^\circ=4,\qquad V=\frac13\cdot4\cdot4=\frac{16}3=5\frac13.\]
\Javob{C) $5\frac13$}

%================= 31
\Savol{31}
$2^6-2^4=64-16=48$ (3 ham, 5 ham yöq qism töplamlar ayirildi).
\Javob{B) 48}

%================= 32
\Savol{32}
$11k\le1000\Rightarrow k\le90$. $11k$ 2 ga ham, 3 ga ham bölinmaydi $\iff$ $k$ bölinmaydi:
\[90-45-30+15=30.\]
\Javob{D) 30}

%================= 33-35
\newpage
\Savol{33}
Hajm $V$. \ Konus: öq kesimi muntazam uçburçak, $h=r_k\sqrt3$, \ $V=\dfrac{\sqrt3}3\pi r_k^3$.\\
Silindr: $h=2r_s$, \ $V=2\pi r_s^3$. \quad Şar: $V=\dfrac43\pi r_{sh}^3$.

\textbf{33.} $2r_s^3=\dfrac{\sqrt3}3r_k^3\Rightarrow\left(\dfrac{r_s}{r_k}\right)^3=\dfrac{\sqrt3}6\Rightarrow\dfrac{r_s}{r_k}=\dfrac{\sqrt[6]3}{\sqrt[3]6}$.

\textbf{34.} $\dfrac43r_{sh}^3=\dfrac{\sqrt3}3r_k^3\Rightarrow\left(\dfrac{r_{sh}}{r_k}\right)^3=\dfrac{\sqrt3}4\Rightarrow\dfrac{r_{sh}}{r_k}=\dfrac{\sqrt[6]3}{\sqrt[3]4}=\dfrac{\sqrt[6]3\cdot\sqrt[3]2}{2}=\dfrac{\sqrt[6]{12}}2$.

\textbf{35.} $2r_s^3=\dfrac43r_{sh}^3\Rightarrow\left(\dfrac{r_s}{r_{sh}}\right)^3=\dfrac23\Rightarrow\dfrac{r_s}{r_{sh}}=\sqrt[3]{\dfrac23}=\dfrac{\sqrt[3]{18}}3$.
\Javob{33 – B; \ 34 – A; \ 35 – E}

%================= 36
\newpage
\Savol{36}
$x^3-1=(x-1)(x^2+x+1)$. \ $u=x^2+x+1$, \ $v=x-1$:
\[2u^2-13uv-7v^2=0\ \Rightarrow\ (2u+v)(u-7v)=0.\]
$2u+v=0$: $2x^2+3x+1=0\Rightarrow x=-1;\ -\frac12$. \quad $u=7v$: $x^2-6x+8=0\Rightarrow x=2;\ 4$.

a) 4 ta ildiz.

b) $(-1)\cdot\left(-\frac12\right)\cdot2\cdot4=4$.
\Javob{a) 4; \quad b) 4}

%================= 37
\newpage
\Savol{37}
$(\sin x\pm\cos x)^2=1\pm\sin2x$:
\[(1+\sin2x)^2+(1-\sin2x)^2+\sin4x=3\ \Rightarrow\ 2\sin^22x+2\sin2x\cos2x=1=\sin^22x+\cos^22x.\]
\[\sin^22x+2\sin2x\cos2x-\cos^22x=0.\]
$\cos2x=0$ bölsa $\sin^22x=0$ — mumkin emas. $\cos^22x$ ga bölamiz: $\tg^22x+2\tg2x-1=0\Rightarrow\tg2x=-1\pm\sqrt2$.\\
$\tg2x=\sqrt2-1=\tg\frac\pi8\Rightarrow x=\dfrac\pi{16}+\dfrac{\pi k}2$; \quad $\tg2x=-1-\sqrt2=\tg\left(-\frac{3\pi}8\right)\Rightarrow x=-\dfrac{3\pi}{16}+\dfrac{\pi k}2$.

a) Eng kiçik musbat ildiz: $\dfrac\pi{16}$.

b) $[-\pi;2\pi]$ (uzunligi $3\pi$, har oilaning davri $\frac\pi2$): \ birinçi oila $k=-2,\dots,3$ — 6 ta; ikkinçi oila $k=-1,\dots,4$ — 6 ta. Jami 12.
\Javob{a) $\frac\pi{16}$; \quad b) 12}

%================= 38
\newpage
\Savol{38}
$x\to1-x$:
\[\begin{cases}2f(x)+f(1-x)=\dfrac{3x^2+x+4}{x(1-x)},\\[6pt] f(x)+2f(1-x)=\dfrac{3x^2-7x+8}{x(1-x)}.\end{cases}\]
Birinçisini 2 ga köpaytirib, ikkinçisini ayiramiz:
\[3f(x)=\frac{3x^2+9x}{x(1-x)}\ \Rightarrow\ f(x)=\frac{x+3}{1-x}.\]
a) $f(2)=\dfrac5{-1}=-5$.

b) $f(x)=-2$: $x+3=-2+2x\Rightarrow x=5$, ya'ni $f^{-1}(-2)=5$.
\Javob{a) $-5$; \quad b) 5}

%================= 39
\newpage
\Savol{39}
a) $f'(x)=\dfrac1{\sqrt x}=\tg30^\circ=\dfrac1{\sqrt3}\Rightarrow x_0=3$, \ $f(3)=2\sqrt3$. Urinma: $y=2\sqrt3+\dfrac{x-3}{\sqrt3}$, \ $y=0\Rightarrow x=-3$.\\
$P(-3;0)$ dan grafikdagi $(x;2\sqrt x)$ nuqtagaça: $d^2=(x+3)^2+4x=x^2+10x+9$ — $x\ge0$ da ösuvçi, eng kiçigi $x=0$ da:
\[d_{\min}=\sqrt9=3.\]
b) $x-2y+6=0$; \ grafik nuqtasi $(t^2;2t)$:
\[d=\frac{|t^2-4t+6|}{\sqrt5}=\frac{(t-2)^2+2}{\sqrt5}\ge\frac2{\sqrt5}=\frac{2\sqrt5}5\quad(t=2).\]
\Javob{a) 3; \quad b) $\frac{2\sqrt5}5$}

%================= 40
\newpage
\Savol{40}
$-1$ — oddiy ildiz, $2$ — karrali ildiz (uriniş): $f(x)=(x+1)(x-2)^2=x^3-3x^2+4$.

a) $a+b+c=f(1)-1=2-1=1$.

b) $t=x-2$:
\[S=\int_{-1}^2(x+1)(x-2)^2dx=\int_{-3}^0(t+3)t^2\,dt=\left(\frac{t^4}4+t^3\right)\Big|_{-3}^0=-\left(\frac{81}4-27\right)=\frac{27}4.\]
\Javob{a) 1; \quad b) $\frac{27}4$}

%================= 41
\newpage
\Savol{41}
$AD=2$, $AN=2+4=6$, $AL=12$, $KL=6$. $AK$ töğri çiziq: $\triangle ARD\sim\triangle ATN\sim\triangle AKL$:
\[RD=\frac{AD\cdot KL}{AL}=\frac{2\cdot6}{12}=1,\qquad TN=\frac{6\cdot6}{12}=3.\]
a) $DRTN$ — trapetsiya: $S=\dfrac{1+3}2\cdot4=8$.

b) $TK$: gorizontal $NL=6$, vertikal $KL-TN=3$ (\textcolor{aux}{töğri burçakli uçburçak}): \ $TK=\sqrt{36+9}=3\sqrt5$.
\Javob{a) 8; \quad b) $3\sqrt5$}

%================= 42
\newpage
\Savol{42}
$BC=5$, \ $\cos\angle ACB=\frac35$.

a) $\angle MCF=360^\circ-90^\circ-90^\circ-\angle ACB=180^\circ-\angle ACB$, \ $\cos\angle MCF=-\frac35$. $\triangle MCF$ da kosinuslar teoremasi:
\[MF^2=3^2+5^2+2\cdot3\cdot5\cdot\frac35=52\ \Rightarrow\ MF=2\sqrt{13}.\]
b) $\angle BAC=\angle BO_2C=90^\circ$ $\Rightarrow$ $A,B,O_2,C$ bir aylanada; $O_2B=O_2C$ $\Rightarrow$ $AO_2$ — $\angle A$ bissektrisasi: $\angle O_2AC=45^\circ$.\\
Ptolemey: $AO_2\cdot BC=AB\cdot O_2C+AC\cdot O_2B\Rightarrow AO_2\cdot5=7\cdot\dfrac5{\sqrt2}\Rightarrow AO_2=\dfrac7{\sqrt2}$.\\
$AO_1$ — kvadrat diagonalining yarmi: $AO_1=\dfrac3{\sqrt2}$, \ $\angle O_1AC=45^\circ$ (pastga). Demak $\angle O_1AO_2=90^\circ$:
\[S=\frac12\cdot\frac3{\sqrt2}\cdot\frac7{\sqrt2}=\frac{21}4.\]
\Javob{a) $2\sqrt{13}$; \quad b) $\frac{21}4$}

%================= 43
\newpage
\Savol{43}
Aylanaga içki çizilgan trapetsiya teng yonli: $AB=CD$, $AC=BD=10$. $B$ dan $AD$ ga balandlik asosi $A$ dan $\frac{13-3}2=5$, $D$ dan $8$ uzoqlikda: $h=\sqrt{100-64}=6$, \ $AB=\sqrt{25+36}=\sqrt{61}$.

a) $BE=BD$ $\Rightarrow$ $\smile BAE=\smile BCD$, ya'ni $\smile BA+\smile AE=\smile BC+\smile CD$. $\smile BA=\smile CD$ $\Rightarrow$ $\smile AE=\smile BC$ $\Rightarrow$ $AE=BC=3$.

b) $ABDE$ uçun Ptolemey: $BE\cdot AD=BD\cdot AE+DE\cdot AB$:
\[10\cdot13=10\cdot3+DE\cdot\sqrt{61}\ \Rightarrow\ DE=\frac{100}{\sqrt{61}}.\]
$\angle AED=180^\circ-\angle ABD$; \ $S_{ABD}=\frac12\cdot13\cdot6=39\Rightarrow\sin\angle ABD=\dfrac{2\cdot39}{10\sqrt{61}}=\dfrac{39}{5\sqrt{61}}$.
\[S_{AED}=\frac12\cdot3\cdot\frac{100}{\sqrt{61}}\cdot\frac{39}{5\sqrt{61}}=\frac{1170}{61},\qquad S_{ABCD}=\frac{3+13}2\cdot6=48.\]
\[S_{ABCDE}=48+\frac{1170}{61}=\frac{4098}{61}.\]
\Javob{a) 3; \quad b) $\frac{4098}{61}$}

%================= 44
\newpage
\Savol{44}
$r=3$; konus balandligi $\sqrt{5^2-3^2}=4$; silindr uzunligi $15-3-4=8$.

a) $S=\underbrace{2\pi\cdot9}_{\text{yarim şar}}+\underbrace{2\pi\cdot3\cdot8}_{\text{silindr yoni}}+\underbrace{\pi\cdot3\cdot5}_{\text{konus yoni}}=18\pi+48\pi+15\pi=81\pi$.

b) $V=\frac23\pi\cdot27+\pi\cdot9\cdot8+\frac13\pi\cdot9\cdot4=18\pi+72\pi+12\pi=102\pi$.
\Javob{a) $81\pi$; \quad b) $102\pi$}

%================= 45
\newpage
\Savol{45}
a) Uç $O(0;0)$ da: $y=kx^2$. Katta ustun uçi $(150;67)$: $k=\dfrac{67}{22\,500}$, \ $y=\dfrac{67}{22\,500}x^2$.

b) Har tomonda 6 ta kiçik ustun: $x=150-15m$, $m=1,\dots,6$, ya'ni $x=135,120,105,90,75,60$.
\[2k\left(135^2+120^2+105^2+90^2+75^2+60^2\right)=2\cdot\frac{67}{22\,500}\cdot225\cdot(81+64+49+36+25+16)=\frac{67\cdot271}{50}=363{,}14.\]
\Javob{a) $y=\frac{67}{22500}x^2$; \quad b) 363,14}

\ShowJavoblar
\end{document}
